\chapter{Método de análisis de las implementaciones}

Esta parte estudia cómo los principios de COBRIA aparecen en dos sistemas existentes sin
convertir el manuscrito en documentación de productos. Los casos se mantienen anonimizados
como Sistema A y Sistema B. El primero está orientado a eventos, secuencias, contenido
multimodal y componentes de inteligencia; el segundo es una plataforma empresarial
multiinquilino con un dominio amplio, operaciones críticas y una cantidad elevada de
entidades.

El análisis se realizó sobre una instantánea del código disponible el 20 de agosto de
2026. Los conteos representan archivos estructurales encontrados en esa fecha y no
usuarios, transacciones, cobertura funcional ni calidad. Los nombres de clases,
organizaciones, personas, rutas físicas, credenciales y reglas comerciales particulares
no se publican.

\section{Preguntas del estudio de implementación}

La observación busca responder cinco preguntas:

\begin{enumerate}[leftmargin=*]
  \item ¿Cómo se distribuyen modelo, acceso a datos, servicios, casos de uso, esquemas y
  catálogos?
  \item ¿Qué elementos corresponden a patrones conocidos y cuáles constituyen la
  integración COBRIA?
  \item ¿Dónde aparecen entidades canónicas y proyecciones o índices derivados?
  \item ¿Qué tan consistente es la aplicación de la estructura dentro de cada caso?
  \item ¿Qué cambios serían necesarios para implementar COBRIA de manera verificable y
  no solo nominal?
\end{enumerate}

\section{Fuentes de evidencia}

Se utilizaron archivos de objetos, modelos base, Daters, conectores, servicios, casos de
uso, esquemas, registries, catálogos, documentación técnica y pruebas existentes. La
presencia de un archivo indica intención estructural; solo el código ejecutable y las
pruebas pueden apoyar una afirmación de comportamiento.

\begin{table}[H]
\centering
\caption{Niveles de evidencia utilizados.}
\label{tab:niveles-evidencia-implementacion}
\begin{tabularx}{\textwidth}{p{2cm}p{4cm}X}
\toprule
Nivel & Evidencia & Interpretación permitida \\
\midrule
E0 & Nombre o documento & Intención; no prueba implementación.\\
E1 & Clase o contrato & Capacidad modelada.\\
E2 & Flujo conectado & Capacidad integrada en código.\\
E3 & Prueba automatizada & Comportamiento verificado en entorno controlado.\\
E4 & Telemetría reproducible & Comportamiento observado bajo carga declarada.\\
\bottomrule
\end{tabularx}
\end{table}

Este criterio evita atribuir funcionalidad completa a una clase porque su nombre
contenga términos como inteligencia, sincronización o auditoría. También impide usar el
número de archivos como sustituto de mantenibilidad.

\section{Procedimiento de inventario}

El procedimiento fue: identificar directorios de capas; clasificar archivos por función;
examinar bases y ejemplos representativos; buscar relaciones entre objeto y Dater;
revisar conectores y rutas; localizar índices, cachés, metadatos y servicios; y contrastar
el hallazgo con la documentación. Los resultados se sintetizaron en categorías neutrales.

Los conteos se calcularon mediante un criterio reproducible: archivos JavaScript bajo
los directorios de objetos, Daters, servicios, casos de uso, schemas y catálogos o datos
de referencia. Archivos generados y manuales pueden incrementar estas cifras; por eso se
reportan como tamaño estructural aproximado.

\section{Neutralidad y ética}

No se copiaron datos operacionales ni secretos. Los ejemplos de esta parte son
pseudocódigo o estructuras reducidas. Las descripciones de dominio se mantienen al nivel
necesario para explicar variación arquitectónica. Un lector puede reproducir la
implementación de referencia sin conocer la identidad de los sistemas originales.

\chapter{Implementación neutral de referencia}

\section{Objetivo}

La implementación neutral no busca reemplazar las aplicaciones estudiadas. Su función es
mostrar la unidad mínima de COBRIA sin depender de un marco de trabajo o proveedor. Debe ser
pequeña, ejecutable y suficientemente explícita para probar identidad, ámbito, mapeo,
proyección, reconstrucción y autorización.

\section{Estructura propuesta}

\begin{verbatim}
src/
  core/
    entities/          entidades canónicas
    value-objects/     valores inmutables
    policies/          invariantes de dominio
  application/
    use-cases/         operaciones autorizadas
    ports/             contratos de persistencia y eventos
  data/
    mappers/           dominio <-> representación física
    repositories/      acceso por identidad y consultas declaradas
    projections/       contratos, writers y rebuilders
    connectors/        adaptadores del proveedor
  metadata/
    schemas/           tipos y validación
    catalogs/          vocabularios y referencias
    relationships/     vínculos semánticos
  analytics/
    conjuntos de datos/          contratos y manifiestos
    características/          definiciones fuera de línea/en línea
    lineage/           grafo de procedencia
  intelligence/
    retrieval/         fragmentos, embeddings y búsqueda
    models/            registro y evaluación
    tools/             capacidades gobernadas para agentes
  governance/
    policies/          acceso, retención y clasificación
    evidence/          auditoría y evidencia
\end{verbatim}

La estructura expresa responsabilidades, no una obligación de crear un archivo por cada
elemento. Un sistema pequeño puede agrupar módulos siempre que mantenga las fronteras.

\section{Entidad canónica}

Una implementación mínima contiene identidad, ámbito, revisión, versión de esquema,
estado y atributos. El constructor valida forma básica; las transiciones con significado
de negocio se realizan mediante métodos o políticas.

\begin{verbatim}
class CanonicalEntity {
  constructor({ id, ámbito, revision = 0, schemaVersion = 1,
                status = "active", attributes = {} }) {
    requireId(id); requireScope(ámbito);
    this.id = id;
    this.ámbito = ámbito;
    this.revision = revision;
    this.schemaVersion = schemaVersion;
    this.status = status;
    this.attributes = structuredClone(attributes);
  }
}
\end{verbatim}

El objeto no conoce URL, SDK, tabla, colección ni caché. Esta independencia corresponde
al principio de separar dominio y persistencia y se relaciona con Mapeador de datos y
Repositorio \parencite{fowler2002poeaa}.

\section{Mapeador y Repositorio}

El Mapeador define correspondencia reversible entre nombres de dominio y campos físicos.
El Repositorio ofrece operaciones en lenguaje del dominio y exige ámbito en todas ellas.

\begin{verbatim}
class EntityMapper {
  toDomain(id, ámbito, crudo) { /* mapeo, valores predeterminados y versión */ }
  toRaw(entity)            { /* representación persistible */ }
}

class EntityRepositorio {
  get(ámbito, id)                 { /* lectura por clave */ }
  save(ámbito, entity, revision)  { /* control optimista */ }
  listBy(ámbito, query, cursor)   { /* consulta permitida */ }
}
\end{verbatim}

El Repositorio no debe devolver el SDK del proveedor ni permitir que la presentación
forme rutas. Un Conector traduce sus operaciones a la API concreta.

\section{Dater como variante de acceso a datos}

Los dos casos utilizan el término local \emph{Dater} para clases que convierten objetos,
ejecutan CRUD, manejan caché y, en algunos módulos, mantienen índices. El nombre no
designa un patrón reconocido distinto: funcionalmente combina responsabilidades de Data
Acceso Objeto, Repositorio, Mapeador de datos e Mapa de identidad. COBRIA puede conservar el término
por compatibilidad, pero recomienda declarar cuál de esas responsabilidades cumple cada
clase.

\begin{table}[H]
\centering
\caption{Correspondencia del Dater con patrones consolidados.}
\label{tab:dater-patrones}
\begin{tabularx}{\textwidth}{p{3.2cm}p{4cm}X}
\toprule
Responsabilidad observada & Patrón relacionado & Recomendación \\
\midrule
CRUD del proveedor & Data Acceso Objeto & Mantener detrás de un contrato.\\
Obtener entidad por identidad & Repositorio & Usar vocabulario del dominio.\\
Convertir crudo y objeto & Mapeador de datos & Separar mapeo cuando crece.\\
Conservar instancia por ID & Mapa de identidad & Definir ámbito e invalidación.\\
Listas con TTL & Caché & Medir frescura y fallos.\\
Índices derivados & Vista materializada & Versionar y reconstruir.\\
\bottomrule
\end{tabularx}
\end{table}

\section{Contrato de proyección ejecutable}

\begin{verbatim}
const ProjectionContract = {
  id: "entity-by-status",
  version: 2,
  sourceType: "Entidad",
  dependencies: [],
  consistency: "eventual",
  transform(entity) { return [[entity.status, entity.id, true]]; },
  verify(source, projection) { /* cobertura y huérfanos */ },
  rebuild({ ámbito, cursor })  { /* lotes y checkpoints */ }
};
\end{verbatim}

El contrato reúne propiedades que en las implementaciones aparecen distribuidas entre
rutas, métodos de Dater, procesos de trabajo y convenciones. Su versión permite crear una proyección
paralela y efectuar cambio de versión.

\section{Caso de uso y transacción lógica}

Un caso de uso coordina autorización, carga, invariantes, persistencia, cola de salida y resultado.

\begin{verbatim}
UpdateEntity(command, actor):
  policy.require(actor, "entity.update", command.ámbito)
  current <- repository.require(command.ámbito, command.id)
  updated <- domain.apply(current, command.change)
  repository.save(command.ámbito, updated, current.revision)
  cola de salida.append(EntityChanged(updated.id, updated.revision))
  evidence.record(command, actor, updated)
  return presenter.summary(updated)
\end{verbatim}

Si el motor permite actualización multipath atómica, canónico, índices críticos y
cola de salida pueden compartir una frontera. Si no, se documenta la compensación y se emplean
procesos de trabajo idempotentes.

\section{Pruebas mínimas}

La referencia requiere pruebas de ida y vuelta del Mapeador, aislamiento de ámbito,
revisión optimista, idempotencia, reconstrucción desde cero, reconstrucción incremental,
detección de huérfanos, cambio de versión y rechazo de escritura directa en una
proyección. Para analítica agrega corte temporal, linaje y equivalencia fuera de línea--en línea.

\chapter{Caso de estudio A: sistema orientado a eventos y contenido}

\section{Caracterización estructural}

En la fecha de corte, el inventario del Sistema A encontró aproximadamente 88 archivos
de objetos de dominio, 91 Daters, 182 servicios, siete casos de uso explícitos, dos
schemas centrales y 102 archivos de catálogos o datos de referencia. Estas cantidades
describen una base heterogénea: módulos antiguos conviven con subsistemas recientes más
formalizados.

El dominio combina participantes, grupos, sesiones, eventos, acciones evaluadas,
planificación, salud, inventario, archivos multimedia, facturación y artefactos de IA.
La variedad permite observar si una estructura común se adapta a información temporal,
conectada y multimodal.

\begin{table}[H]
\centering
\caption{Inventario neutral del Sistema A.}
\label{tab:inventario-sistema-a}
\begin{tabularx}{\textwidth}{p{3.5cm}p{2.2cm}X}
\toprule
Categoría & Conteo & Observación \\
\midrule
Objetos de dominio & 88 & Estilos de modelado de distintas generaciones.\\
Daters & 91 & CRUD, caché, serialización e índices.\\
Servicios & 182 & Operación, analítica, multimedia e IA.\\
Casos de uso explícitos & 7 & Concentrados en flujos recientes.\\
Esquemas centrales & 2 & La validación también aparece distribuida.\\
Catálogos y referencias & 102 & Vocabularios, configuración y datos compartidos.\\
\bottomrule
\end{tabularx}
\end{table}

\section{Objetos y serialización}

Los objetos observados usan clases JavaScript con constructores, getters, setters y, en
varios módulos, métodos de serialización. Algunos emplean campos privados y copias
defensivas para estructuras anidadas; otros exponen propiedades públicas y conservan
alias para compatibilidad. Esta coexistencia demuestra adaptabilidad, pero también riesgo
de semántica desigual.

En los módulos de inteligencia, los objetos representan versiones de modelos, métricas,
experimentos, consentimientos, tareas, auditorías, conocimiento, retroalimentación y solicitudes
de acción. Esto se alinea con la Parte IV porque el modelo y la decisión no se reducen a
un valor suelto: poseen identidad, estado, procedencia y tiempos.

\section{Daters e índices derivados}

Los Daters del Sistema A acceden al conector de datos, convierten cargas útiles, mantienen
Mapa locales y actualizan índices por modelo, tipo, estado u otras dimensiones. Algunos
construyen una actualización multipath que persiste el registro principal y varias rutas
derivadas. Este es un ejemplo concreto del intercambio de COBRIA: una escritura lógica
produce varias escrituras físicas para acelerar lecturas.

\begin{figure}[H]
\centering
\begin{tikzpicture}[
  node distance=.8cm and .85cm,
  box/.style={draw,align=center,minimum width=2.3cm,minimum height=.82cm},
  derived/.style={box,fill=cobriaLight},
  line/.style={-{Latex[length=2mm]},thick}
]
\node[box] (object) {Objeto de\\dominio};
\node[box,right=of object] (dater) {Dater};
\node[box,right=of dater] (connector) {Conector\\NoSQL};
\node[derived,below=of connector] (base) {Ruta\\principal};
\node[derived,left=of base] (indexes) {Índices por\\dimensión};
\node[derived,left=of indexes] (cache) {Mapa e índices\\locales};
\draw[line] (object) -- (dater);
\draw[line] (dater) -- (connector);
\draw[line] (connector) -- (base);
\draw[line] (connector) -- (indexes);
\draw[line] (dater) -- (cache);
\end{tikzpicture}
\caption{Flujo de datos observado y abstraído en el Sistema A.}
\label{fig:flujo-sistema-a}
\end{figure}

La ruta principal se aproxima al canónico y los índices a proyecciones, pero esa
clasificación debe verificarse por entidad. Un carga útil duplicado puede contener campos
que otro módulo modifica directamente; en tal caso existe autoridad compartida y no una
proyección COBRIA válida.

\section{Servicios, casos de uso e IA}

La capa de servicios incluye coordinación de eventos, análisis, reportes, sincronización,
multimedia, salud, facturación e inteligencia. Los casos de uso explícitos aparecen en
operaciones que requieren transiciones y permisos. En otros módulos, un servicio o Dater
todavía concentra la coordinación.

La capacidad de IA es estructuralmente relevante: existen artefactos para versiones,
experimentos, métricas, preferencias, consentimiento, conocimiento, auditoría, retroalimentación
y acciones propuestas. Esto apoya la adaptabilidad conceptual de COBRIA, pero no prueba
exactitud de modelos ni seguridad de agentes. Esas capacidades deben alcanzar E3 o E4
mediante los experimentos de la Parte IV.

\section{Fortalezas observadas}

\begin{itemize}[leftmargin=*]
  \item separación visible entre presentación, lógica y datos;
  \item objetos especializados para eventos, multimedia e inteligencia;
  \item Daters con cachés e índices adecuados para consultas frecuentes;
  \item actualizaciones multipath en operaciones con derivados;
  \item servicios para analítica y ciclo de vida de artefactos;
  \item adopción progresiva de casos de uso en flujos con reglas.
\end{itemize}

\section{Deuda y oportunidades}

Los principales riesgos son estilos de objetos heterogéneos, validación distribuida,
responsabilidades amplias en ciertos Daters, índices sin contrato uniforme, rutas
conocidas por múltiples módulos y cobertura desigual de casos de uso. La optimización
local mediante Mapa puede perder ámbito o frescura si la clave e invalidación no se
declaran.

La migración recomendada no reescribe todos los módulos. Se priorizan entidades con
datos sensibles, múltiples índices, IA o alto volumen. Para cada una se registra fuente
canónica, contrato de proyección, versión, propietario, reconstrucción y pruebas. Los
alias de compatibilidad pueden permanecer en el Mapeador mientras se estabiliza el modelo.

\section{Evolución longitudinal posterior del Sistema A}

Una segunda observación del repositorio, posterior al inventario inicial, permitió
analizar cambio arquitectónico y no solo estructura estática. La evolución no se usa
como un nuevo prueba comparativa de rendimiento; constituye evidencia de aplicabilidad de los
principios de COBRIA en una migración real y gradual.

\begin{table}[H]
\centering
\caption{Evidencia longitudinal neutral del Sistema A.}
\label{tab:evolucion-sistema-a}
\begin{tabularx}{\textwidth}{p{4.1cm}p{2.3cm}X}
\toprule
Propiedad observada & Evidencia & Interpretación COBRIA \\
\midrule
Frontera Presentación--Data & 0 importaciones directos & La dependencia se hace verificable
por CI, no solo convencional.\\
Compatibilidad temporal & 134 consumidores inventariados & La deuda se congela y solo
puede disminuir.\\
Vocabulario físico central & 36 nodos semánticos & Los nombres compactos quedan detrás
de un registro de rutas.\\
Modelos de lectura explícitos & 4 repositorios de vistas & Listados y resúmenes declaran
almacén derivado separado.\\
Pruebas arquitectónicas y de escala & 16 artefactos relevantes & Se verifican rutas,
concurrencia, vistas, presupuestos y fronteras.\\
\bottomrule
\end{tabularx}
\end{table}

Las nuevas vistas de listas y resúmenes incluyen \texttt{complete},
\texttt{updatedAt}, TTL y operaciones de reemplazo, invalidación o reconstrucción.
Una vista fría agrupa solicitudes concurrentes en una sola reconstrucción; las pruebas
comprueban que veinte solicitudes comparten un único warmup. La autorización se ejecuta
en cada solicitud contra la política canónica, aun cuando el resultado materializado
esté fresco. Este detalle separa optimización de autoridad.

Las operaciones de dominio aplican actualizaciones multipath para canónico, índices,
auditoría y notificación; utilizan versiones, transiciones explícitas, idempotencia y
controles concurrentes para invariantes. Además, el contexto utilizado por telemetría y
FinOps excluye valores no validados provenientes del cliente. Estas decisiones amplían
COBRIA con dos reglas: \emph{la optimización puede fallar sin falsificar el canónico} y
\emph{la atribución operacional debe proceder del ámbito autorizado}.

La migración del interfaz aporta otra lección. Una frontera temporal puede conservar
134 consumidores heredado sin permitir nuevos consumidores. El manifiesto y el límite
automático convierten el patrón estrangulador en una transición medible. COBRIA debe, por
tanto, exigir a cada puente: propietario, inventario, fecha o condición de retiro,
presupuesto que no crezca y pruebas de delegación.

\section{Implicación para la validez del caso}

La observación longitudinal fortalece H8 porque el mismo sistema evolucionó hacia rutas
semánticas, servicios de aplicación, políticas separadas, proyecciones materializadas y
controles automatizados arquitectónicos sin una reescritura total. No demuestra causalidad: los cambios
fueron realizados por participantes conocedores de la propuesta y no se compararon
contra un equipo de control. Sí proporciona evidencia concreta de que COBRIA puede
funcionar como arquitectura evolutiva y no únicamente como taxonomía retrospectiva.

\chapter{Caso de estudio B: sistema empresarial multiinquilino}

\section{Caracterización estructural}

El Sistema B posee una escala estructural considerablemente mayor. En la fecha de corte
se encontraron aproximadamente 680 objetos de dominio, 680 Daters, 1,624 servicios, 18
casos de uso explícitos, 686 schemas y 514 archivos de catálogos o datos de referencia.
La correspondencia cercana entre objetos, Daters y schemas revela una estrategia
sistemática, parte de ella apoyada por generación o convenciones repetibles.

\begin{table}[H]
\centering
\caption{Inventario neutral del Sistema B.}
\label{tab:inventario-sistema-b}
\begin{tabularx}{\textwidth}{p{3.5cm}p{2.2cm}X}
\toprule
Categoría & Conteo & Observación \\
\midrule
Objetos de dominio & 680 & Cobertura de numerosos subdominios.\\
Daters & 680 & Correspondencia casi uno a uno con objetos.\\
Servicios & 1,624 & Políticas, procesos de trabajo, contratos y orquestadores.\\
Casos de uso explícitos & 18 & Operaciones críticas seleccionadas.\\
Esquemas & 686 & Validación estructural extensa.\\
Catálogos y referencias & 514 & Vocabularios y configuración por dominio.\\
\bottomrule
\end{tabularx}
\end{table}

Estos números no permiten concluir que todo módulo esté integrado o probado. Sí ofrecen
evidencia de una arquitectura dirigida por estructura y metadatos, adecuada para estudiar
escalabilidad organizacional.

\section{Modelo base y mapeo físico}

Un modelo base representativo inicializa identidad y valores predeterminados, restringe campos
desconocidos, serializa a JSON y mantiene un mapa explícito entre campo físico y campo de
dominio. El método inverso reconstruye una instancia desde datos crudo. Esta forma
materializa Mapeador de datos de manera directa.

La ventaja es que una entidad puede conservar vocabulario estable aunque cambien claves
físicas. El riesgo aparece cuando valores predeterminados silenciosos ocultan ausencia o versiones
antiguas. COBRIA recomienda distinguir campo ausente, nulo y vacío y registrar la versión
del esquema utilizada en la reconstrucción.

\section{Dater base}

El Dater base observado recibe ruta, modelo, conector y TTL. Implementa lectura por ID,
alta, actualización, borrado, lista, paginación por cursor, consulta por hijo, carga
múltiple, escucha en tiempo real, caché de entidades y caché de listas. Convierte entre
crudo y objeto mediante el modelo.

\begin{figure}[H]
\centering
\begin{tikzpicture}[
  node distance=.75cm and .8cm,
  box/.style={draw,align=center,minimum width=2.25cm,minimum height=.82cm},
  derived/.style={box,fill=cobriaLight},
  line/.style={-{Latex[length=2mm]},thick}
]
\node[box] (service) {Servicio /\\caso de uso};
\node[box,right=of service] (dater) {Dater base\\especializado};
\node[box,right=of dater] (mapper) {Modelo y\\mapa de campos};
\node[box,below=of mapper] (connector) {Conector};
\node[derived,left=of connector] (store) {Almacenamiento\\multiinquilino};
\node[derived,left=of store] (cache) {Caché con\\TTL};
\draw[line] (service) -- (dater);
\draw[line] (dater) -- (mapper);
\draw[line] (dater) -- (connector);
\draw[line] (connector) -- (store);
\draw[line] (dater) -- (cache);
\end{tikzpicture}
\caption{Abstracción del acceso común observado en el Sistema B.}
\label{fig:flujo-sistema-b}
\end{figure}

Esta base reduce duplicación y ofrece una API uniforme. Sin embargo, el argumento de
ruta debe incorporar ámbito de forma no opcional. Una ruta base global o una clave de
caché basada solo en ID podría romper aislamiento aunque el resto de la aplicación
conozca el organización.

\section{Metadatos, relaciones y catálogos}

El Sistema B contiene registries para políticas de campo, relaciones, clasificaciones de
referencias y correspondencia con catálogos. Los campos pueden declararse configurables,
externos o asociados con vocabularios. Esta capa convierte reglas que normalmente
quedarían dispersas en datos inspeccionables.

La arquitectura dirigida por metadatos favorece generación de objetos, schemas, Daters,
documentación y controles de interfaz. Para que la generación sea confiable debe existir
una fuente primaria del metamodelo, validación del generador y detección de divergencia.
Editar artefactos generados manualmente sin retroalimentar el metamodelo crea deriva.

\section{Servicios y operaciones críticas}

Los servicios abarcan consulta, disponibilidad, inventario, finanzas, comunicaciones,
integraciones, privacidad, sincronización, reportes y operaciones multiinquilino. Existen
contratos de repositorio, procesos de trabajo, políticas, state machines y orquestadores. Los casos de
uso explícitos se concentran en operaciones donde estados, permisos y efectos derivados
requieren una frontera clara.

Los contratos operacionales muestran ámbito de organización y unidad, claves de
idempotencia, autenticación, recálculo en servidor, rate limits y errores semánticos.
Estos mecanismos coinciden con COBRIA Núcleo y Gobierno. La persistencia de cola de salida,
operaciones idempotentes, versiones de esquema, bloqueos y trabajos programados ofrece
puntos naturales para reconstrucción y observabilidad.

\section{Fortalezas observadas}

\begin{itemize}[leftmargin=*]
  \item correspondencia sistemática entre objeto, Dater y esquema;
  \item modelo base con mapeo entre dominio y representación física;
  \item Dater base con caché, paginación y tiempo real;
  \item metadatos explícitos para campos, relaciones y catálogos;
  \item contratos de repositorio y casos de uso en operaciones críticas;
  \item idempotencia, cola de salida, procesos de trabajo y controles multiinquilino modelados;
  \item suficiente variedad de dominios para evaluar extensibilidad.
\end{itemize}

\section{Deuda y oportunidades}

La escala produce riesgos de proliferación: una clase por estructura puede aumentar
costo de navegación, generación y pruebas. Una correspondencia uno a uno no garantiza
que cada estructura posea identidad o comportamiento suficiente para ser entidad. Los
1,624 servicios pueden incluir políticas pequeñas y artefactos generados; deben
agruparse por responsabilidad y dependencia antes de inferir complejidad.

COBRIA recomienda clasificar cada tipo como entidad, value object, evento, comando,
proyección o catálogo. También propone registrar responsabilidad de proyecciones, reconstrucción
y presupuesto. La base común debe exigir ámbito, revisión y versión; los servicios deben
depender de contratos, no de rutas ni SDKs.

\chapter{Comparación entre casos y derivación de COBRIA}

\section{Similitudes estructurales}

Ambos sistemas separan presentación, lógica y datos; representan el dominio mediante
objetos; encapsulan persistencia en Daters; utilizan servicios; conservan catálogos y
datos compartidos; y desnormalizan para responder consultas. Ambos poseen conectores a
un almacenamiento NoSQL y cachés locales. Estas similitudes motivaron la búsqueda de un
vocabulario común.

\begin{longtable}{>{\bfseries}p{3.2cm}p{5.6cm}p{5.6cm}}
\caption{Comparación neutral de los casos.}\label{tab:comparacion-casos}\\
\toprule
Dimensión & Sistema A & Sistema B \\
\midrule
\endfirsthead
\toprule
Dimensión & Sistema A & Sistema B \\
\midrule
\endhead
Dominio & Eventos, secuencias, multimedia e inteligencia. & Operaciones empresariales
multiinquilino de gran amplitud.\\
Objetos & Heterogéneos y especializados. & Sistemáticos, con modelo base frecuente.\\
Daters & Especializados; varios mantienen índices propios. & Base reutilizable y amplia
correspondencia por entidad.\\
Esquemas & Parciales o distribuidos. & Cobertura estructural extensa.\\
Metadatos & Catálogos y configuraciones por módulo. & Registries de campos, relaciones y
catálogos.\\
Casos de uso & Adopción focalizada en flujos recientes. & Adopción focalizada en
operaciones críticas.\\
Proyecciones & Índices locales/remotos y artefactos analíticos. & Índices, reportes,
cachés, cola de salida y snapshots.\\
IA & Componentes explícitos de versión, auditoría y retroalimentación. & Preparación estructural,
automatización y datos empresariales.\\
Riesgo principal & Heterogeneidad y contratos implícitos. & Proliferación y divergencia
entre artefactos.\\
\bottomrule
\end{longtable}

\section{Diferencias que prueban adaptabilidad}

El Sistema A necesita representar tiempo fino, acciones, secuencias y archivos. El
Sistema B necesita ámbito jerárquico, catálogos configurables, operaciones idempotentes y
muchas relaciones. COBRIA no intenta igualar sus modelos. La adaptación ocurre porque
ambos pueden expresar identidad, ámbito, versión, mapeo y proyección, pero emplean
contratos diferentes.

En A, una proyección puede indexar eventos por participante, contexto o ventana y un
pipeline puede generar clips o características. En B, una proyección puede resolver
disponibilidad, saldos, reportes o búsquedas por organización. El mecanismo es común; la
semántica pertenece al dominio.

\section{Elementos heredados y contribución}

Los objetos, DAO/Repositorio, Mapeador de datos, Capa de servicios, casos de uso, caché, catálogos,
vistas materializadas, cola de salida y metadatos son conceptos existentes. El Dater es un nombre
local para una combinación de esos patrones, no una nueva categoría científica.

La contribución COBRIA se ubica en la integración de estas ideas alrededor de:

\begin{enumerate}[leftmargin=*]
  \item una fuente canónica con identidad, ámbito, revisión y versión;
  \item una separación explícita entre dominio y forma física;
  \item proyecciones con contrato, responsabilidad, consistencia y reconstrucción;
  \item metadatos que conectan modelo, esquema, relaciones y catálogos;
  \item continuidad de linaje desde operación hasta conjunto de datos, modelo y decisión;
  \item una frontera que impide a derivados o agentes adquirir autoridad por accidente.
\end{enumerate}

\section{Patrones de diseño identificados}

\begin{table}[H]
\centering
\caption{Patrones observados y papel dentro de COBRIA.}
\label{tab:patrones-observados}
\begin{tabularx}{\textwidth}{p{3.4cm}p{4.2cm}X}
\toprule
Patrón & Evidencia típica & Papel COBRIA \\
\midrule
Repositorio/DAO & Daters y contratos de repositorio & Aislar persistencia.\\
Mapeador de datos & fromRaw, toRaw y mapas de campos & Separar semántica y forma física.\\
Mapa de identidad & Mapas por identidad & Evitar instancias duplicadas en sesión.\\
Capa de servicios & Servicios de dominio y aplicación & Coordinar operaciones.\\
Caso de uso & Flujos explícitos y autorizados & Definir frontera transaccional.\\
Vista materializada & Índices, resúmenes y snapshots & Optimizar lectura con reconstrucción.\\
CQRS parcial & Modelos de consulta especializados & Separar necesidades de lectura.\\
Cola de salida & Eventos pendientes y procesos de trabajo & Reducir ventana de pérdida.\\
Registro & Metadatos y catálogos & Centralizar contratos inspeccionables.\\
Estrategia y política & Validadores y políticas intercambiables & Variar reglas sin rutas físicas.\\
Adaptador/Pasarela & Conectores y proveedores externos & Encapsular infraestructura.\\
Máquina de estados & Transiciones de operaciones críticas & Restringir ciclos de vida.\\
\bottomrule
\end{tabularx}
\end{table}

\section{Modelo de madurez}

\begin{table}[H]
\centering
\caption{Niveles de madurez COBRIA para una entidad.}
\label{tab:madurez-cobria}
\begin{tabularx}{\textwidth}{p{1.5cm}p{3.2cm}X}
\toprule
Nivel & Nombre & Evidencia \\
\midrule
M0 & Acceso directo & interfaz de usuario o servicio conoce rutas y cargas útiles.\\
M1 & Encapsulado & Objeto y Dater/Repositorio básico.\\
M2 & Canónico & Identidad, ámbito, mapeo y esquema explícitos.\\
M3 & Reconstruible & Proyecciones versionadas con rebuild y verificación.\\
M4 & Gobernado & Linaje, políticas, observabilidad y recuperación.\\
M5 & Adaptativo & Optimización basada en evidencia con aprobación.\\
\bottomrule
\end{tabularx}
\end{table}

La madurez se evalúa por entidad o flujo, no por sistema completo. Ambos casos pueden
contener simultáneamente módulos M1 y M4. Esta granularidad permite una migración basada
en riesgo.

\chapter{Estrategia de adopción y migración}

\section{Adopción incremental}

La estrategia recomendada sigue siete pasos:

\begin{enumerate}[leftmargin=*]
  \item seleccionar una entidad con dolor verificable;
  \item declarar identidad, ámbito, autoridad e invariantes;
  \item aislar Mapeador y Repositorio o documentar el Dater existente;
  \item inventariar copias, índices, cachés y consumidores;
  \item crear contratos de proyección y verificación;
  \item instrumentar lecturas, escrituras, anomalías y costo;
  \item reconstruir en paralelo, comparar y efectuar cambio de versión reversible.
\end{enumerate}

No se recomienda renombrar masivamente los Daters. La migración comienza por contratos y
propiedades verificables. El nombre puede cambiar cuando aporte claridad y no rompa
consumidores.

\section{estrangulador para datos}

Una fachada compatible intercepta operaciones antiguas y las dirige al nuevo Repositorio.
Durante la transición puede comparar lecturas o emitir eventos para nuevas proyecciones.
La escritura dual sin verificación aumenta riesgo; se prefiere una fuente autoritativa y
un derivado reconstruible.

\begin{figure}[H]
\centering
\begin{tikzpicture}[
  node distance=.85cm and .9cm,
  box/.style={draw,align=center,minimum width=2.35cm,minimum height=.82cm},
  derived/.style={box,fill=cobriaLight},
  line/.style={-{Latex[length=2mm]},thick}
]
\node[box] (heredado) {Consumidor\\existente};
\node[box,right=of heredado] (facade) {Fachada\\compatible};
\node[derived,right=of facade] (usecase) {Nuevo caso\\de uso};
\node[derived,below=of usecase] (repo) {Repositorio y\\Mapeador};
\node[box,left=of repo] (canonical) {Canónico};
\node[derived,left=of canonical] (projection) {Proyección\\versionada};
\draw[line] (heredado) -- (facade);
\draw[line] (facade) -- (usecase);
\draw[line] (usecase) -- (repo);
\draw[line] (repo) -- (canonical);
\draw[line] (canonical) -- (projection);
\end{tikzpicture}
\caption{Migración incremental mediante una fachada compatible.}
\label{fig:migracion-strangler}
\end{figure}

\section{Priorización por riesgo y valor}

Se asigna una puntuación basada en sensibilidad, frecuencia, amplificación, número de
copias, incidentes, dificultad de reconstrucción y consumidores. Las entidades con alto
riesgo y alto uso se atienden primero. Un catálogo estático y estable puede permanecer
en M1 o M2 si no existe beneficio neto de mayor formalidad.

\section{Generación dirigida por metadatos}

En el Sistema B, la escala hace atractiva la generación. El metamodelo puede producir
esqueleto de entidad, esquema, Mapeador, Dater/Repositorio, documentación y pruebas de
contrato. La generación debe ser determinista; el repositorio rechaza divergencias entre
salida generada y archivos versionados.

El generador no decide invariantes de negocio ni crea proyecciones por cada campo. Esas
decisiones requieren demanda de consulta, responsable y presupuesto.

\section{Compatibilidad y versionado}

Los cambios se clasifican como aditivos, deprecatorios o incompatibles. Un Mapeador puede
leer varias versiones y escribir solo la vigente. La migración se ejecuta por lotes y
con checkpoint. Las proyecciones incluyen versión en ruta o clave para permitir
coexistencia y reversión.

\section{Presupuesto de deuda y retiro de puentes}

Toda compatibilidad temporal se registra en un manifiesto reproducible. Sea \(L_t\) el
conjunto de consumidores heredado permitidos en la versión \(t\). La regla de integración
continua exige:

\[
L_{t+1}\subseteq L_t.
\]

Un consumidor nuevo requiere migración al contrato vigente, no ampliación silenciosa
del manifiesto. El puente solo reexporta o adapta; no recibe nuevas reglas de negocio. Su
eliminación se realiza por dominio, con pruebas de comportamiento y reducción del
presupuesto después de cada bloque estable.

\section{Criterio de finalización}

Una entidad se considera migrada cuando: toda escritura autorizada atraviesa casos de
uso definidos; la fuente canónica es única; el ámbito es obligatorio; el ida y vuelta del
Mapeador está probado; cada derivado posee contrato y responsable; una reconstrucción completa
produce el resultado esperado; la telemetría permite medir costo y anomalías; y los
controles automatizados impiden reintroducir rutas físicas, dependencias invertidas o consumidores heredado.

\chapter{Reproducibilidad, limitaciones y síntesis de la Parte V}

\section{Paquete reproducible previsto}

El manuscrito debe acompañarse de un paquete que no contenga los proyectos privados. El
paquete incluirá implementación neutral, generadores sintéticos, contratos de ejemplo,
scripts de inventario, pruebas, configuración del prueba comparativa y plantillas de resultados.
Cada versión publicada fijará versiones y resúmenes criptográficos.

\begin{table}[H]
\centering
\caption{Contenido del paquete reproducible.}
\label{tab:paquete-reproducible}
\begin{tabularx}{\textwidth}{p{3.6cm}X}
\toprule
Artefacto & Contenido \\
\midrule
Referencia & Núcleo, Repositorio, Mapeador, Gestor de proyecciones y Gobierno.\\
Datos sintéticos & Perfiles A y B con semilla y escalas configurables.\\
Contratos & Entidades, proyecciones, conjuntos de datos y herramientas.\\
Migraciones & Total, incremental, cambio de versión y reversión.\\
Pruebas & Aislamiento, idempotencia, ida y vuelta y reconstrucción.\\
Prueba comparativa & Cargas, instrumentación y análisis estadístico.\\
Documentación & Diccionario, decisiones y guía de réplica.\\
\bottomrule
\end{tabularx}
\end{table}

\section{Limitaciones del análisis estático}

El inventario no demuestra qué archivos se ejecutan en producción, qué rutas contienen
datos, cuántos módulos tienen consumidores ni qué cobertura alcanzan las pruebas. Los
conteos pueden incluir archivos generados, contratos pequeños o capacidades parciales.
La documentación puede estar desactualizada. Estas limitaciones se resolverán en la
Parte VI mediante ejecución controlada, trazas y experimentos.

\section{Riesgo de sesgo del investigador}

El investigador es también autor o conocedor de los sistemas, lo que facilita acceso y
comprensión, pero puede favorecer interpretaciones positivas. Se mitigará publicando
criterios previos, scripts, resultados negativos, niveles de evidencia y una
implementación neutral. Cuando sea posible, otro evaluador revisará la clasificación de
una muestra de módulos.

\section{Validez de la generalización}

Los casos comparten lenguaje y almacenamiento NoSQL, por lo que no prueban adaptación a
todos los ecosistemas. La contribución generalizable es el mecanismo y sus condiciones,
no los nombres de carpetas ni las cantidades. Una réplica futura debe incluir otro
lenguaje, otro motor documental y un equipo independiente.

\section{Síntesis de la Parte V}

Esta parte conectó la formalización de COBRIA con evidencia estructural de dos sistemas
reales y distintos. El Sistema A mostró objetos heterogéneos, Daters con índices y una
amplia superficie de eventos, multimedia e inteligencia. El Sistema B mostró una
correspondencia masiva entre objetos, Daters y schemas, un modelo base con mapeo físico,
metadatos y operaciones multiinquilino más sistemáticas.

El análisis confirmó que ``Objeto--Dater--Servicio--Catálogo'' no constituye por sí solo
un patrón nuevo. Es una composición local de Modelo de dominio, DAO/Repositorio, Mapeador de datos,
Capa de servicios, Registro, Caché y vistas materializadas. COBRIA adquiere interés de
investigación al convertir esa composición en una arquitectura explícita, formal y
evaluada alrededor de autoridad canónica, proyecciones reconstruibles, ámbito y linaje
hacia IA.

También se definió una implementación neutral, un modelo de madurez y una estrategia de
migración incremental. Los conteos estructurales sirven para caracterizar, no para
celebrar tamaño. La siguiente parte ejecutará el protocolo: medirá rendimiento,
reconstrucción, aislamiento, evolución, linaje y adaptabilidad, y presentará tanto
resultados favorables como negativos.
